\documentclass[11pt]{article}

\usepackage[margin=1in]{geometry}
\usepackage{enumitem}
\usepackage{amsmath, amssymb}
\usepackage{hyperref}

\setlist[enumerate]{itemsep=0.8em, topsep=0.8em}
\setlist[itemize]{itemsep=0.3em, topsep=0.3em}

\begin{document}

\begin{center}
  {\Large \textbf{CHEM 350 Examination}}\\[0.25em]
  {\large \textbf{Organic Chemistry: Reactions and Mechanisms}}\\[0.5em]
\end{center}

\noindent\textbf{Instructions:}
\begin{itemize}
  \item Answer all questions.
  \item For Questions 1--4, choose the best option.
  \item For Questions 5--8, write brief but precise answers.
\end{itemize}

\vspace{0.5em}

\begin{enumerate}[label=\textbf{\arabic*.}]

  \item Which factor primarily determines whether an $\mathrm{S_N1}$ or $\mathrm{S_N2}$ mechanism 
  will predominate in a nucleophilic substitution reaction?
  \begin{enumerate}[label=(\Alph*)]
    \item The strength of the nucleophile only
    \item The structure of the alkyl halide (primary, secondary, tertiary)
    \item The temperature of the reaction
    \item The polarity of the solvent exclusively
  \end{enumerate}

  \item In electrophilic aromatic substitution, which substituent on a benzene ring is both 
  ortho/para-directing and deactivating?
  \begin{enumerate}[label=(\Alph*)]
    \item Methyl group ($-\mathrm{CH_3}$)
    \item Nitro group ($-\mathrm{NO_2}$)
    \item Halogen atoms (e.g., $-\mathrm{Cl}$, $-\mathrm{Br}$)
    \item Amino group ($-\mathrm{NH_2}$)
  \end{enumerate}

  \item What is the primary product when 1-butene undergoes hydrohalogenation with HBr in the 
  presence of peroxides (anti-Markovnikov conditions)?
  \begin{enumerate}[label=(\Alph*)]
    \item 1-bromobutane
    \item 2-bromobutane
    \item 1,2-dibromobutane
    \item 2-bromo-2-methylpropane
  \end{enumerate}

  \item Which spectroscopic technique is most useful for distinguishing between constitutional 
  isomers based on their connectivity?
  \begin{enumerate}[label=(\Alph*)]
    \item UV-Vis spectroscopy
    \item Infrared (IR) spectroscopy
    \item Nuclear Magnetic Resonance (NMR) spectroscopy
    \item Mass spectrometry
  \end{enumerate}

  \item Explain the concept of carbocation stability and how it influences the regioselectivity 
  of addition reactions to alkenes. Use Markovnikov's rule as an example.

  \item Describe the mechanism of the aldol condensation reaction. What conditions favor the 
  formation of the $\alpha,\beta$-unsaturated carbonyl product rather than just the aldol adduct?

  \item Explain the difference between kinetic and thermodynamic control in organic reactions. 
  Provide an example where reaction conditions can be tuned to favor one product over another.

  \item Discuss the role of resonance stabilization in determining the acidity of organic compounds. 
  Compare the acidity of phenol, ethanol, and acetic acid, explaining the structural features that 
  account for their differences.

\end{enumerate}

\end{document}
